% --- Document: Preamble ---
\documentclass[../main.tex]{subfiles}
\graphicspath{{\subfix{../assets/}}}
% --- Document: Start ---
\begin{document}
% Introduction
\section{Fundamentação Teórica}    
% Project Problens 
\subsection{Gestão imobiliária de Contratos de Locação}
DESCREVER OS PROBLEMAS QUE VISAMOS RESOLVER, COMO CONTROLE DE CONTRATOS, GERAÇÃO DE DESPESAS E DESCONTRO, CONTROLE DE CONTAS E FORMAS DE pagamento
\subsubsection{Tipos de Locações}
Segundo Barbosa (2020), “a legislação brasileira vigente estabelece cinco categorias principais de locações urbanas: residencial, não residencial, 
por temporada, mista e institucional, cada uma com características jurídicas próprias e regulação específica na Lei do Inquilinato”. A locação residencial é 
voltada à moradia do locatário e de sua família, sendo regulada pelos artigos 47 a 50 da Lei nº 8.245/1991, que tratam das hipóteses de retomada do imóvel e prorrogação contratual. 
A locação não residencial, por sua vez, destina-se ao exercício de atividades comerciais, industriais ou profissionais, conforme estabelecido nos artigos 51 a 57 da mesma lei, 
que garantem, por exemplo, o direito à renovação compulsória do contrato desde que cumpridos certos requisitos. A locação por temporada é disciplinada nos artigos 48 a 50 e 
caracteriza-se pelo uso transitório do imóvel por prazo não superior a 90 dias, sendo frequentemente utilizada para fins de lazer, trabalho temporário, cursos ou tratamento de saúde. 
A locação mista, embora não prevista expressamente em capítulo específico da Lei do Inquilinato, é reconhecida pela doutrina e pela jurisprudência como aquela em que o imóvel é utilizado 
tanto para fins residenciais quanto comerciais, aplicando-se a ela as disposições gerais da legislação conforme o uso predominante. Por fim, a locação institucional, também não regulamentada de 
forma detalhada na Lei nº 8.245/1991, ocorre em situações especiais envolvendo entes públicos ou instituições privadas para instalação de serviços ou atividades de interesse público, sendo comumente 
regida por normas específicas ou contratos administrativos, mas ainda sujeita, de forma subsidiária, aos princípios da legislação locatícia vigente.
\subsubsection{Sistema imobiliário de Locação}

\subsubsection{Gestão imobiliária}
descricao do problema 2
\subsubsection{Legislação imobiliária}
descricao do problema 2
\subsubsection{Obrigações Legais e Judiciais}
descricao do problema 2

% Project Organization
\subsection{Administração de Caixa}
DESCRECER A PROPOSTA A PARA A SOLUÇÃO DOS PROBLEMAS ENCONTRADOS
\subsubsection{Administração de Valores a Receber}
descricao da solucao do problema 1
\subsubsection{Administração de Valores a Pagar}
descricao da solucao do problema 2
\subsubsection{Administração de Capital de Giro}

% --- Document: End ---
\end{document}

